% !PRINTED 364
\weberpage{364}

Um hieraus die Beziehung zwischen $h$ und $h'$ herzuleiten, machen wir von den allgemeinen Gruppensätzen Bd.~II, \S\,2 Gebrauch, die wir folgendermaßen formulieren und ergänzen:

12. \gesperrt{Satz. Ist $A$ eine Gruppe und $B$ ein Teiler von $A$ von endlichem Index $(A,B)$, ferner $C$ ein Teiler von $B$ von endlichem Index $(B,C)$, so ist Bd.~II, \S\,2, (4):}
\weq{4}{(A,C) = (A,B)(B,C).}
Ist $A$ wieder eine Gruppe mit dem Teiler $B$ vom Index $\mu$, so setzen wir
\weq{5}{A = a_1 B + a_2 B + \cdots + a_\mu B,}
worin $a_1, a_2, \dots, a_\mu$ ein volles Repräsentantensystem von $A$ nach $B$ ist.

Es sei nun $C$ eine Gruppe von Elementen, die mit den Elementen von $A$ zusammensetzbar sind (z.\,B. in unserem Falle einfach durch Multiplikation), so ergibt sich aus (5):
\weq{6}{AC = a_1 BC + a_2 BC + \cdots + a_u BC.} % [sic] Gedruckt ist u statt \mu.
Nun kann in den beiden Nebengruppen $a_1BC$, $a_2BC$ nur dann ein und dasselbe Element vorkommen, wenn $a_1a_2^{-1}$ in $BC$, aber nicht in $B$ enthalten ist.

Daraus formulieren wir den Satz:

13. \gesperrt{Satz. Ist $A$ eine Gruppe, $B$ ein Teiler von $A$ vom endlichen Index $(A,B)$, ferner $C$ eine mit $A$ zusammensetzbare Gruppe von der Art, daß $A$ und $BC$ außer den $B$ keine gemeinschaftlichen Elemente enthalten, so ist}
\weq{7}{(AC,\ BC) = (A,B).}
Die Voraussetzung dieses Satzes läßt sich auch so aussprechen:

\begin{center}\gesperrt{$B$ ist der Durchschnitt von $A$ und $BC$,}\end{center}

und sie ist z.\,B. erfüllt, wenn $C$ in $B$ enthalten, also $B = BC$ ist, und auch dann, wenn $C$ mit $A$ kein Element gemein hat.

14. \gesperrt{Satz. Ist $A$ eine aus den beiden Gruppen $A'$ und $B$ zusammengesetzte Gruppe}
\weq{8}{A = A'B}
\gesperrt{und $B'$ der Durchschnitt von $A'$ und $B$, so ist}
\weq{9}{(A,B) = (A',B'),}
\gesperrt{vorausgesetzt, daß diese Indices endlich sind.}

Denn nach (8) ist in jeder Nebengruppe $aB$ zu $B$ in $A$ ein Element $a'$ aus $A'$ enthalten, und man kann also diese Neben-
%CONT>>
